\documentclass{article}
\usepackage[a4paper, margin=1in]{geometry}
\usepackage{datetime}
\usepackage{wrapfig}
\usepackage[pdftex]{graphicx} \pdfcompresslevel=9
\newdate{date}{21}{09}{2026}
\date{21/09/2026}
\usepackage[french]{babel}
\usepackage{url}

\title{TELECOM Graphisme par ordinateur: Introduction shaders}
\author{camille.schreck@inria.fr}

\begin{document}

\maketitle

\url{https://www.shadertoy.com}

\begin{enumerate}
\item Faire un dégradé de gauche à droite.\\
  {\centering
    \includegraphics[width=0.3\linewidth]{greyscale.png}
  }
\item Trouver une colormap, puis l'appliquer au dégradé.\\
  {\centering
    \includegraphics[width=0.3\linewidth]{magma.png}
  }
\item Faire en sorte de pouvoir choisir une direction.\\
  {\centering
    \includegraphics[width=0.3\linewidth]{magma_r.png}
  }
\item Dessiner un cercle au centre.
\item Créer une fonction \texttt{bool circle(vec2 p, float r, vec2 c)} qui renvoit true si p est dans le cercle de rayon r et de centre c. L'utiliser pour dessiner plusieurs cercle.
\item Utiliser la variable \texttt{iTime} de Shadertoy pour faire bouger les cercles.
\item Créer un système planétaire (soleil, planetes, satellites).\\
  {\centering
    \includegraphics[width=0.3\linewidth]{circles.png}
  }
\item Faites une scène plus complexe.\\ 
  {\centering
    \includegraphics[width=0.3\linewidth]{planetes.png}
  }
\end{enumerate}

\end{document}
